% !PRINTED 602
\weberpage{602}

In $Q$ kommt also ein dreigliedriger Cyklus vor.

2. Es enthalte $\varkappa$ zwei dreigliedrige Cyklen $(1,\,2,\,3)(4,\,5,\,6)$. Wir nehmen $\pi = (1,\,3,\,4)$ an und erhalten

\[
\lambda = (1,\,3,\,2)(4,\,6,\,5)(3,\,2,\,4)(1,\,5,\,6) = (1,\,2,\,5,\,3,\,4).
\]

Diese Permutation $\lambda$, die in $Q$ enthalten ist, fällt aber unter den Fall 1.

3. Es enthalte $\varkappa$ einen dreigliedrigen und einen zweigliedrigen Cyklus $(1,\,2,\,3)(4,\,5)$. (Dass in $\varkappa$, wenn es zu $A$ gehört, noch ein zweiter Cyclus von gerader Gliederzahl vorkommen muss, ist hier gleichgültig). Für $\pi = (1,\,2,\,4)$ ergiebt sich

\[
\lambda = (1,\,3,\,2)(4,\,5)(2,\,4,\,3)(1,\,5) = (1,\,2,\,5,\,3,\,4),
\]

was wieder unter den Fall 1. fällt.

4. Es enthalte $\varkappa$ drei Transpositionen, $(1,\,2)(3,\,4)(5,\,6)$. Für $\pi = (1,\,3,\,5)$ folgt

\[
\lambda = (1,\,2)(3,\,4)(5,\,6)(3,\,2)(5,\,4)(1,\,6) = (1,\,3,\,5)(2,\,6,\,4),
\]

was unter den Fall 2. fällt.

5. Es enthalte $\varkappa$ zwei Transpositionen und ein unverändertes Element $(1,\,2)(3,\,4)(5)$. Man setzt $\pi = (1,\,2,\,5)$ und erhält

\[
\lambda = (1,\,2)(3,\,4)(5)(2,\,5)(3,\,4)(1) = (1,\,5,\,2).
\]

Damit sind alle Fälle erschöpft, wenn $n > 4$ ist. Für $n = 4$ bleibt der eine Fall noch übrig, dass $\varkappa$ aus zwei Transpositionen besteht, der eben das besondere Verhalten bei $n = 4$ herbeiführt, wodurch die algebraische Auflösung der Gleichung $4^{\text{ten}}$ Grades ermöglicht wird\footnote{Der erste vollständige Beweis, dass die allgemeine Gleichung von höherem als dem $4^{\text{ten}}$ Grade durch Radicale nicht lösbar ist, rührt von \gesperrt{Abel} her (Crelle's Journal, Bd I, 1826). Ueber die älteren Beweisversuche von \gesperrt{Ruffini} (1799 bis 1806) und ihr Verhältniss zum \gesperrt{Abel}'schen Beweis vergleiche man die Abhandlung von \gesperrt{Burkhardt}, „Die Anfänge der Gruppentheorie und \gesperrt{Paolo Ruffini}`` (Abhandlungen zur Geschichte der Mathematik VI. Supplement zu Schlömilch's Zeitschrift. Leipzig 1892).}.

Es folgt aus diesem Satze weiter, dass die symmetrische Gruppe keine anderen normalen Theiler hat, als sich selbst, die alternirende Gruppe und die Einheitsgruppe. Denn ist $S$ die symmetrische, $A$ die alternirende Gruppe und $Q$ ein normaler Theiler von $S$, so ist der grösste gemeinschaftliche Theiler $Q'$ von $A$ und $Q$ ein normaler Theiler von $A$, ist also gleich $A$ oder gleich 1.
